\section{The two-channel decomposition and the path of the state}
\label{app:decomposition-path}

This section holds three results that the main text states in words. The forward-looking
gap splits exactly into the target's part and the sender's part
(Section~\ref{sec:results:wedge}); when initial conditions matter, the limiting sanction is
positive (Section~\ref{sec:results:main}); and a sanction rule that does not fall with the
state produces no conflict cycle (Section~\ref{sec:discussion:path}).

\subsection{The path of the state under a sanction rule}
\label{app:state-path}

\paragraph{Dependence on initial conditions implies a positive limiting sanction.}
\begin{proposition}[Limit dependence and the limiting sanction]
\label{prop:limit-dependence}
Let the transition satisfy zero-sanction descent, $\F(x,0)<x$ for every
$x\in(0,1]$ (the erosion half of strict adjustment), and let $R$ be a
continuous feasible sanction rule. If the natural paths of $R$ from two
starts converge to two different limits, then from at least one of them the
limiting sanction $R(x_\infty)$ is strictly positive: a limit state $x_\infty$
of a path under $R$ is a fixed point of $x\mapsto\F(x,R(x))$, and a fixed point
with $R(x_\infty)=0$ can only be $x_\infty=0$
(\path{rulePathF_limit_isFixed}, \path{exists_pos_limit_sanction_of_limit_dependence},
\path{Theory/LockIn/LimitDependence.lean}). The converse fails: at the affine
member of the quadratic--linear family under the lifting condition the non-compliance equilibrium's closed loop has one
globally stable steady state with a positive sanction
(\path{placableAffine_unique_attracting_state}).
\end{proposition}

\paragraph{A non-decreasing rule produces no conflict cycle.}
\begin{proposition}[No conflict cycles under a non-decreasing rule]
\label{prop:no-cycles}
Let $R$ be a sanction rule that maps $[0,1]$ into itself and is
non-decreasing in the state there, and let the state follow
$x_{t+1}=\F(x_t,R(x_t))$ from $x_0\in[0,1]$, with sanctions $s_t=R(x_t)$. Then:
\begin{enumerate}[label=(\roman*)]
\item the state path and the sanction path are monotone in $t$, in the same
direction, and both converge;
\item if $x_{n+k}=x_n$ for some $n\ge0$ and some $k\ge1$, then $x_{n+j}=x_n$ for
every $j\ge0$; in particular a periodic path is constant.
\end{enumerate}
Both conclusions hold in every equilibrium in $\mathcal E$ whose sanction rule is
non-decreasing in the state, and in the constructed equilibrium of each of the four cases of
\Cref{thm:four-cells} at the hypotheses under which it is constructed, that of case
$11$ under the increasing-differences condition \eqref{eq:rally-margin} or with a non-decreasing
response.
\end{proposition}

\begin{proof}
By \textup{(F2)}, \textup{(F3)} and the monotonicity of $R$, the closed loop
$G(y)=\F(y,R(y))$ is non-decreasing on $[0,1]$ (\path{ruleStepF_le_of_le}). If
$x_1\ge x_0$, induction gives $x_{t+1}=G(x_t)\ge G(x_{t-1})=x_t$ for every $t$,
and if $x_1\le x_0$ the reverse; the sanctions $R(x_t)$ move in the same
direction (\path{rulePathF_monotone_or_antitone},
\path{rulePathF_ruleSanctionsF_monotone_or_antitone}). Both sequences are
bounded, so both converge, with no continuity of $R$
(\path{rulePathF_ruleSanctionsF_tendsto}); when $R$ is continuous the limit state
is a rest point of $G$, as in Step~7 of the proof of
\Cref{thm:regime-separation}. For (ii), a monotone path with $x_{n+k}=x_n$ and
$k\ge1$ has $x_{n+1}=x_n$, so $x_n$ is a rest point of $G$ and the path stays
there (\path{rulePathF_const_of_return}, \path{rulePathF_periodic_const}). The
equilibrium form takes $R$ to be the sanction rule of the equilibrium, whose values lie in $[0,1]$
(\path{equilibrium_no_cycles}). In the four cases the rule is zero in $00$ and
$01$ (\path{myopicRule_eq_zero}, \path{senderOnlyRule_eq_zero'}), the myopic
sanction $\sigM(x,0)$ in $10$, non-decreasing by \textup{(S7)}
(\path{targetOnlyRule_monotoneOn}), and the dynamic best response at full non-compliance in $11$,
non-decreasing under the increasing-differences condition (\path{responsePath_monotoneOn_of_margin});
the four statements are \path{myopic_no_cycles}, \path{senderOnly_no_cycles},
\path{targetOnly_no_cycles} and \path{dynamic_no_cycles}
(\path{dynamicOf_no_cycles} with a non-decreasing response). General in $\F$.
\end{proof}

Along deterministic equilibrium paths in this model, a conflict cycle therefore
needs a sanction rule that falls with the state somewhere. No equilibrium in the
paper is shown to have one: the dual-basin equilibrium's rule is shown to be
non-decreasing only on its two basins, and \Cref{prop:identification} does not
cover a falling rule. The cycles of \citet{AcemogluWolitzky2014} rest on
misperceived actions and limited memory, which this model does not have.

\subsection{The forward-looking gap splits into two components}
\label{app:two-channel}

The dynamic--myopic comparison is described in the text as two forces: the target
incumbent's own valuation of the political state its compliance leads to, and the
switch from the sender's myopic best response to its dynamic best response. That
description is exact, and the ordering condition the argument consumes is weaker
than the one stated in \Cref{sec:results}. Both points are recorded here.

\paragraph{Three objectives and two components.}
Fix a feasible state $x$ and a target continuation value $\Vt$, and write, as
functions of compliance $a\in[0,1]$,
\begin{align}
J^{MM}_x(a)&:=\Ut\bigl(x,a,\sigM(x,a)\bigr),
\label{eq:JMM}\\
J^{TM}_x(a)&:=\Ut\bigl(x,a,\sigM(x,a)\bigr)
   +\beta_T\,\Vt\bigl(\F(x,\sigM(x,a))\bigr),
\label{eq:JTM}\\
J^{FF}_x(a)&:=\Ut\bigl(x,a,\sigDynTop_{\Vs}(x,a)\bigr)
   +\beta_T\,\Vt\bigl(\F(x,\sigDynTop_{\Vs}(x,a))\bigr).
\label{eq:JFF}
\end{align}
$J^{MM}$ is the objective the myopic target maximizes and $J^{FF}$ is the
integrand of the target Bellman equation \eqref{eq:T-bellman}, so both are
objects the equilibrium already carries. The middle object $J^{TM}$, a target
that values continuation but faces the \emph{myopic} sender, is a device. It
is not a third equilibrium concept, and no maximizer of it is constructed
anywhere below.

The total forward-looking gap is $\Delta_x:=J^{FF}_x-J^{MM}_x$. Its two
components are the two steps between the three objectives:
\begin{equation}
C_x:=J^{TM}_x-J^{MM}_x,
\qquad
S_x:=J^{FF}_x-J^{TM}_x.
\label{eq:two-channels}
\end{equation}
$C_x$ is the \emph{target-continuation} component (what the incumbent's own
political-survival motive adds, holding the sender's behavior fixed), and
$S_x$ is the \emph{sender-response-switch} component, what changes for the target
purely because the sender it faces is forward-looking.

\begin{proposition}[Exact two-component decomposition]
\label{prop:two-channel}
At every feasible $x$ and every $a\in[0,1]$,
\begin{equation}
\Delta_x(a)=C_x(a)+S_x(a),
\label{eq:decomposition}
\end{equation}
with no residual and no cross term. Moreover
$C_x(a)=\beta_T\,\Vt\bigl(\F(x,\sigM(x,a))\bigr)$, so the target-continuation
component is $\beta_T$ times the target's continuation value at the successor state, which the separated-primitive
route signs directly.
\end{proposition}

\eqref{eq:decomposition} is an identity, not an approximation; the content is
that \eqref{eq:JMM}--\eqref{eq:JFF} are the right three objectives, and that the
sender response is the only thing that moves between $J^{TM}$ and $J^{FF}$.

The split earns its place because the two components have different logical
status. $C_x$ is signed under three conditions stated separately from the
baseline transition primitives: compliance placates, sanctions pass through
strictly to the successor state, and the target's continuation value is
increasing on the reached successors. Only the first is close to baseline; the
transition primitives give weak monotonicity in the sanction rather than strict,
and an increasing target value is not a baseline property at all. $S_x$ is not
signed even under those. Both sender best responses falling in compliance, the dynamic one lying
above the stage one, and the target's payoff rising in the sanction together do
\emph{not} determine whether $S_x$ rises or falls. The component reads the
target's payoff \emph{difference} between the two responses, which turns on the
\emph{gap} between them and not on their level alone, and that gap may widen
with compliance. Any route to the
compliance ordering must therefore either sign $S_x$ separately or dominate it.
The switch from the stage to the dynamic best response of the sender thus has an
\emph{ambiguous} effect on the target's compliance incentive under the baseline;
that is a scope delimiter, not a claim that the intended ordering typically
fails.

\paragraph{Only the net movement is needed.}
Asking each component to behave separately is stronger than the ordering requires.
Because \eqref{eq:decomposition} is exact, the comparison never sees the components
apart: only their sum. The exact condition is that across any rise in
compliance from $a_L$ to $a_H$,
\begin{equation}
S_x(a_H)-S_x(a_L)\;<\;C_x(a_L)-C_x(a_H),
\label{eq:net-channel}
\end{equation}
which permits $S_x$ to \emph{rise} (to work against the ordering), provided
the target-continuation component gives up strictly more. Condition
\eqref{eq:net-channel} is \emph{equivalent} to a strictly decreasing $\Delta_x$,
so it is an accounting restatement of the gap condition rather than a new
assumption: the route was always net, and reading it component by component conceals
the trade-off. \citet{QuahStrulovici2012} is why a component-by-component sign
argument is stronger than a condition on the aggregate (single crossing is not
preserved by addition except under stringent conditions), but
\eqref{eq:net-channel} is not theirs; it is a corollary of the exact
decomposition above.

Two boundaries travel with it, and both are one-directional. The condition is
\emph{cardinal}, a strict inequality between differences, and it is therefore
\emph{strictly stronger} than the ordinal relation \eqref{eq:sc-lower} below:
\eqref{eq:net-channel} implies \eqref{eq:sc-lower}, and the converse fails.
Single crossing survives an arbitrary strictly increasing transformation of
either objective and a difference condition does not, so the two are ordered
rather than incomparable: on $[0,1]$ the pair $f(a)=a^{3}$, $g(a)=a$ stands in
the ordinal relation while $f-g$ rises over $[9/10,1]$. And the sufficient route
in which the switch gap weakly falls while the target gap strictly falls is
strictly stronger again than \eqref{eq:net-channel}, since it forbids exactly the
adverse switch the net condition tolerates. The two-leg \emph{ordinal} route below is a
different sufficient condition again, constraining each leg separately against
the intermediate objective rather than the aggregate; we do not claim an
implication either way between it and \eqref{eq:net-channel}.

\paragraph{The ordering condition is ordinal.}
\Cref{sec:results} states the comparison under a strictly decreasing total
gap. That is sufficient but stronger than necessary. Say that $f$ is
\emph{strictly single-crossing lower} than $g$ on $A$ when, for all $a_L<a_H$ in
$A$,
\begin{equation}
f(a_L)\le f(a_H)
\quad\Longrightarrow\quad
g(a_L)<g(a_H).
\label{eq:sc-lower}
\end{equation}

\begin{proposition}[Ordinal form of the compliance ordering]
\label{prop:ordinal-ordering}
If $J^{FF}_x$ is strictly single-crossing lower than $J^{MM}_x$ on $[0,1]$, then
every maximizer of $J^{FF}_x$ lies weakly below every maximizer of $J^{MM}_x$. A
strictly decreasing $\Delta_x$ implies \eqref{eq:sc-lower}, so the gap
condition is one route to it. Relation \eqref{eq:sc-lower} is transitive:
if $J^{FF}_x$ is strictly single-crossing lower than $J^{TM}_x$, and $J^{TM}_x$
than $J^{MM}_x$, the conclusion follows for the pair $(J^{FF}_x,J^{MM}_x)$.
\end{proposition}

No continuity, concavity, compactness, or uniqueness enters, and the conclusion
is about the argmax \emph{correspondences} rather than about a chosen selection
from either. The transitivity clause is what makes \eqref{eq:decomposition}
usable rather than merely descriptive: the two components can be discharged one at
a time against $J^{TM}_x$, with different tools and at different strengths,
without a maximizer of $J^{TM}_x$ ever being exhibited.

The relations here are, one for one, those of \citet{AmirRietzke2025}: their
Definition~1 is the weak form below, their identity-difference order is the
cardinal gap condition, their Lemma~2(ii) is the mixed composition, and the
auxiliary problem they insert between two objectives is $J^{TM}_x$; we cite
them for the method. Their Theorem~1(i) also supplies a converse on a
one-dimensional choice set: the weak relation holds if and only if the argmax
sets are strong-set ordered for every ordered pair of feasible sets, so the
whole content sits on two-point sets
(\path{sCLowerOn_iff_pairwise_meet_join}); the weak relation is not one
sufficient condition among others but what the strong-set conclusion is.

\paragraph{Only one leg needs to be strict.}
Asking both legs for \eqref{eq:sc-lower} demands the most from the leg that can
supply the least, since the sender-response-switch component's sign is not fixed by
the primitives at all. Write $f\preceq_{SC}g$ for the two-clause \emph{weak}
form: a weak rise in $f$ forces a weak rise in $g$, and a strict rise forces a
strict rise. Then
\[
J^{FF}_x\preceq_{SC}J^{TM}_x,
\qquad
J^{TM}_x\prec_{SC}J^{MM}_x
\]
already gives $J^{FF}_x\prec_{SC}J^{MM}_x$, so the conclusion of
\Cref{prop:ordinal-ordering} still orders \emph{every} pair of maximizers. The
switch component is asked only not to reverse the comparison; the target's own
continuation motive carries the strictness. The mirror composition (strict
switch, weak target) is available and is not the version we recommend.

Two cautions travel with the weak relation. On its own it orders the argmax
\emph{correspondences} in the strong set order, and hence their extremal
elements, but not an arbitrary selection from each: two constant objectives
satisfy it with every action maximizing both. That is precisely why one leg must
remain strict. And the weak relation is not the weakest hypothesis delivering the
conclusion: two argmax sets can be ordered while it fails away from the
maximizers; its value is that it is selection-robust and does not depend on where
the unknown maximizers sit.

\paragraph{The step for the direct effect of the sender's foresight is local.}
The comparison at a fixed target action, $\sigM(x,a)\le\sigDynTop_{\Vs}(x,a)$, does not need
$\Vs$ monotone at every state. It needs only that the sender's continuation term,
read at the current state as a function of the sanction, does not fall over
feasible sanctions; global monotonicity of $\Vs$ is one way of supplying that.
Weaker still, it suffices that the continuation term not prefer the
dynamic best response to the myopic best response wherever the former lies
below the latter: a condition on the two responses alone. Each of these is a weakening of
its predecessor, so nothing that discharged the earlier form is lost.
